\section{Conjuntos de datos y problemas de distribución}

\subsection{Observación tridimensional de la calidad de los datos}

Los conjuntos de datos de Offline RL no son una caja negra; el docente enfatiza: entender el origen de los datos, su cobertura y la distribución de las recompensas es lo que permite guiar de manera efectiva la elección del algoritmo.

\begin{itemize}
    \item \textbf{Cobertura}: si las combinaciones estado-acción son uniformes, si existen regiones dispersas
    \item \textbf{Estructura de la recompensa}: si la recompensa es densa, si tiene ruido, si tiene valores atípicos
    \item \textbf{Diversidad conductual}: si los datos contienen una sola política de comportamiento o una mezcla de varias
\end{itemize}

\begin{importantbox}{Elegir el algoritmo a partir de la calidad de los datos}
Cuando la cobertura de los datos es amplia y la recompensa es limpia, pueden aprovecharse métodos más agresivos de bootstrapping (como IQL); cuando los datos presentan muestras de baja calidad o la recompensa se desvía, conviene priorizar métodos Conservative o el filtrado de muestras basado en la recompensa.
\end{importantbox}

\begin{table}[H]
\centering
\begin{tabular}{p{0.2\textwidth} p{0.25\textwidth} p{0.25\textwidth} p{0.2\textwidth}}
\toprule
Característica de los datos & Descripción esquemática & Riesgo & Estrategia recomendada \\
\midrule
Alta recompensa, baja cobertura & La misma trayectoria se repite & Sobreajuste, falta de generalización & Reforzar la regularization, incorporar aumento de datos \\
Baja recompensa pero diversa & Varias políticas suboptimal & Q overestimation on rare good actions & Filtrar muestras de baja recompensa, usar expectile distortion \\
Comportamiento mixto & Mezcla de varias políticas sin etiquetar & Difícil definir la política de comportamiento & Entrenamiento por segmentos que da mayor peso a behavior cloning \\
\bottomrule
\end{tabular}
\caption{Elegir la estrategia de offline RL según la distribución de los datos}
\end{table}

\subsection{Resumen de la sección}

El éxito de Offline RL depende de una comprensión precisa de la cobertura de los datos, la calidad de la recompensa y la diversidad conductual. Ajustar la estrategia según las características de los datos (filter/IQL/CQL) es más seguro que fijar un algoritmo.

\subsection{Caso: construcción de un pipeline de datos mixtos}

Los equipos de proyecto suelen mezclar simulator logs, human logs y expert trajectories, y vuelven a ponderar las muestras según la recompensa y el importance sampling:

\begin{enumerate}
    \item Primero se calcula el effective sample size de cada trayectoria.
    \item Los datos se dividen en bins según el reward percentile y se les asigna un sampling weight distinto.
    \item Se usa el expectile de IQL para evaluar la distribución de Q en cada bin, y se vuelve a muestrear en ambos extremos para mantener la estabilidad.
    \item Si el expectile se sale del intervalo previsto, se regresa al modo conservative o se filtra el bin anómalo.
\end{enumerate}

\begin{knowledgebox}{Ventajas del pipeline de datos mixtos}
El sesgo de las distintas fuentes de datos difiere: el simulator es más preciso pero carece de diversidad, mientras que los human data tienen ruido pero son más naturales. El pipeline mixto, mediante re-weighting, permite que el algoritmo sea a la vez robusto y capaz de generalizar.
\end{knowledgebox}
